\subsection{Vacuum-energy correction from the quartic interactions}
\label{subsec:vacuum-energy-correction}

We now consider the vacuum-energy correction generated by the four-derivative
interactions in Eq.~\eqref{eq:two-string-ng-sum-relative}.  The calculation
below makes one additional assumption that is not contained in the
Nambu--Goto Lagrangian itself: the sum modes $X_{\mathrm{s}}^i$ are kept
massless, while the relative modes $X_{\mathrm{r}}^i$ are assigned a common
mass $m$ by adding
\begin{equation}
  \Delta\mathcal L_m
  = -\frac{T}{4}m^2 X_{\mathrm r}^i X_{\mathrm r}^i .
  \label{eq:relative-mass-term}
\end{equation}
The normalization in Eq.~\eqref{eq:relative-mass-term} follows from the
unnormalized sum/relative variables in Eq.~\eqref{eq:sum-relative-fields}.
Because Eq.~\eqref{eq:two-string-ng-sum-relative} itself has an independent
shift symmetry for $X_{\mathrm r}^i$, the origin of the mass term must be
specified by additional physics.  The result below should therefore be read
as the vacuum diagram generated by the Nambu--Goto derivative interactions in
the presence of the quadratic mass term \eqref{eq:relative-mass-term}; a
microscopic mechanism for the mass may generate additional local interactions.

Let $N$ denote the number of massive relative components.  If every transverse
relative mode is massive, then $N=D-2$.  It is convenient to introduce
canonically normalized fields
\begin{equation}
  \phi_{\mathrm s}^i = \sqrt{\frac{T}{2}}\,X_{\mathrm s}^i,
  \qquad
  \phi_{\mathrm r}^i = \sqrt{\frac{T}{2}}\,X_{\mathrm r}^i .
  \label{eq:canonical-sr-fields}
\end{equation}
The quadratic Minkowski Lagrangian is then
\begin{equation}
  \mathcal L_2
  = -\frac12 \partial_a\phi_{\mathrm s}^i\partial^a\phi_{\mathrm s}^i
    -\frac12 \partial_a\phi_{\mathrm r}^i\partial^a\phi_{\mathrm r}^i
    -\frac12 m^2\phi_{\mathrm r}^i\phi_{\mathrm r}^i .
  \label{eq:quadratic-canonical-massive-relative}
\end{equation}
After Wick rotation, the massive propagator is
\begin{equation}
  \langle \phi_{\mathrm r}^i(p)\phi_{\mathrm r}^j(-p)\rangle_0
  = \frac{\delta^{ij}}{p^2+m^2},
  \label{eq:massive-relative-propagator}
\end{equation}
whereas the sum-field propagator is $\delta^{ij}/p^2$.

For reference, define the Euclidean invariants
\begin{align}
  A_{\mathrm s} &\equiv
    \partial_a\phi_{\mathrm s}^i\partial_a\phi_{\mathrm s}^i,
  &
  A_{\mathrm r} &\equiv
    \partial_a\phi_{\mathrm r}^i\partial_a\phi_{\mathrm r}^i,
  \notag\\
  C &\equiv
    \partial_a\phi_{\mathrm s}^i\partial_a\phi_{\mathrm r}^i,
  &
  H^Y_{ab} &\equiv
    \partial_a\phi_Y^i\partial_b\phi_Y^i,
  \qquad Y=\mathrm s,\mathrm r,
  \notag\\
  H^{\mathrm{sr}}_{ab} &\equiv
    \partial_a\phi_{\mathrm s}^i\partial_b\phi_{\mathrm r}^i
   +\partial_a\phi_{\mathrm r}^i\partial_b\phi_{\mathrm s}^i .
\end{align}
The quartic Euclidean interaction obtained from
Eq.~\eqref{eq:two-string-ng-sum-relative} is
\begin{align}
  \mathcal L_{4,E}
  ={}& \frac{1}{16T}(A_{\mathrm s}+A_{\mathrm r})^2
      +\frac{1}{4T}C^2
      -\frac{1}{8T}
       (H^{\mathrm s}_{ab}+H^{\mathrm r}_{ab})
       (H^{\mathrm s}_{ab}+H^{\mathrm r}_{ab})
      -\frac{1}{8T}H^{\mathrm{sr}}_{ab}H^{\mathrm{sr}}_{ab} .
  \label{eq:quartic-euclidean-sr}
\end{align}
The signs in Eq.~\eqref{eq:quartic-euclidean-sr} are worth checking: the
Lorentz-scalar derivative contractions become positive Euclidean contractions,
while $\mathcal L_E=-\mathcal L_M$ after Wick rotation.

At first order in the quartic interaction, the connected contribution to the
Euclidean vacuum functional is a single quartic vertex whose four legs are
contracted pairwise.  Diagrammatically this is the two-loop ``figure-eight''
vacuum graph.  Writing $W=-\log Z$ and dividing by the Euclidean worldsheet
volume gives
\begin{equation}
  \Delta\rho_{\mathrm{fig8}}
  = \left\langle \mathcal L_{4,E}\right\rangle_0 .
  \label{eq:vacuum-energy-first-interaction}
\end{equation}
This is the first correction generated by the quartic interactions.  It should
not be confused with the Gaussian one-loop determinant, which exists already
in the free massive theory.

We use dimensional regularization in $d=2-2\epsilon$ dimensions and an
infinite flat worldsheet.  Define
\begin{equation}
  K_0 \equiv \int_p \frac{p^2}{p^2},
  \qquad
  K_m \equiv \int_p \frac{p^2}{p^2+m^2},
  \qquad
  \int_p \equiv
  \left(\frac{e^{\gamma_E}\bar\mu^2}{4\pi}\right)^\epsilon
  \int\frac{d^d p}{(2\pi)^d} .
  \label{eq:K0-Km-definitions}
\end{equation}
In dimensional regularization $K_0=\int_p 1=0$ because it is scaleless.
Consequently every pure-$\phi_{\mathrm s}$ vacuum contraction is proportional
to $K_0^2$ and every mixed $\phi_{\mathrm s}$--$\phi_{\mathrm r}$ contraction
is proportional to $K_0K_m$; all of them vanish.  Thus, under these specific
infinite-volume dimensional-regularization assumptions, only the purely
massive part contributes,
\begin{equation}
  \mathcal L_{4,E}^{(r)}
  = \frac{1}{16T}
      \left(\partial_a\phi_{\mathrm r}^i
            \partial_a\phi_{\mathrm r}^i\right)^2
    -\frac{1}{8T}
      \left(\partial_a\phi_{\mathrm r}^i
            \partial_b\phi_{\mathrm r}^i\right)
      \left(\partial_a\phi_{\mathrm r}^j
            \partial_b\phi_{\mathrm r}^j\right).
  \label{eq:massive-only-quartic-euclidean}
\end{equation}
This simplification is regulator and geometry dependent.  In particular, on a
finite cylinder the massless sums are not scaleless, so the mixed terms must
be retained when computing a finite-size/Casimir energy.

Rotational and $O(N)$ invariance give the coincident derivative contraction
\begin{equation}
  \left\langle
    \partial_a\phi_{\mathrm r}^i
    \partial_b\phi_{\mathrm r}^j
  \right\rangle_0
  = \delta^{ij}\delta_{ab}\frac{K_m}{d} .
  \label{eq:derivative-tadpole-tensor}
\end{equation}
Introduce
\begin{equation}
  A \equiv
  \partial_a\phi_{\mathrm r}^i\partial_a\phi_{\mathrm r}^i,
  \qquad
  B \equiv
  (\partial_a\phi_{\mathrm r}^i\partial_b\phi_{\mathrm r}^i)
  (\partial_a\phi_{\mathrm r}^j\partial_b\phi_{\mathrm r}^j).
\end{equation}
For $\langle A^2\rangle_0$, the disconnected Wick pairing inside each copy of
$A$ gives $N^2K_m^2$, while the two crossed pairings together give
$2NK_m^2/d$.  Therefore
\begin{equation}
  \langle A^2\rangle_0
  = \frac{N(Nd+2)}{d}\,K_m^2 .
  \label{eq:A2-wick-result}
\end{equation}
For $\langle B\rangle_0$, the three Wick pairings give respectively
$N^2K_m^2/d$, $NK_m^2$, and $NK_m^2/d$, so that
\begin{equation}
  \langle B\rangle_0
  = \frac{N(N+d+1)}{d}\,K_m^2 .
  \label{eq:B-wick-result}
\end{equation}
Substituting Eqs.~\eqref{eq:A2-wick-result} and \eqref{eq:B-wick-result} into
Eq.~\eqref{eq:vacuum-energy-first-interaction} gives the regulated result
\begin{align}
  \Delta\rho_{\mathrm{fig8}}
  &= \frac{1}{16T}\langle A^2\rangle_0
     -\frac{1}{8T}\langle B\rangle_0
  \notag\\
  &= \frac{N}{16Td}
     \left[N(d-2)-2d\right]K_m^2 .
  \label{eq:figure-eight-general-d}
\end{align}
This form should be kept until after renormalization.  Although its strict
$d=2$ coefficient is
\begin{equation}
  \left.\frac{N}{16Td}
  \left[N(d-2)-2d\right]\right|_{d=2}
  = -\frac{N}{8T},
  \label{eq:strict-two-dimensional-coefficient}
\end{equation}
the term proportional to $N(d-2)$ is evanescent and multiplies the divergent
quantity $K_m^2$.  Setting $d=2$ before expanding the regulated integral would
therefore lose part of the single-pole and finite structure.

To evaluate the loop integral, write
\begin{equation}
  K_m
  = \int_p\left(1-\frac{m^2}{p^2+m^2}\right)
  = -m^2 I_m,
  \qquad
  I_m \equiv \int_p\frac{1}{p^2+m^2},
  \label{eq:Km-Im-relation}
\end{equation}
where the first term vanishes as a scaleless integral.  With the
$\overline{\mathrm{MS}}$ measure used in Eq.~\eqref{eq:K0-Km-definitions},
\begin{align}
  I_m
  &= \frac{1}{4\pi}
     e^{\gamma_E\epsilon}\Gamma(\epsilon)
     \left(\frac{\bar\mu^2}{m^2}\right)^\epsilon
  \notag\\
  &= \frac{1}{4\pi}
     \left[
       \frac{1}{\epsilon}
       +L
       +\epsilon\left(\frac{L^2}{2}+\frac{\pi^2}{12}\right)
       +\mathcal O(\epsilon^2)
     \right],
  \qquad
  L\equiv\log\frac{\bar\mu^2}{m^2} .
  \label{eq:massive-tadpole-dimreg}
\end{align}
Hence the compact regulated expression is
\begin{equation}
  \boxed{
  \Delta\rho_{\mathrm{fig8}}
  = \frac{N}{16Td}
    \left[N(d-2)-2d\right]
    m^4 I_m^2
  }
  \qquad (d=2-2\epsilon).
  \label{eq:figure-eight-compact-result}
\end{equation}
For checking pole terms, the coefficient in Eq.~\eqref{eq:figure-eight-general-d}
expands as
\begin{equation}
  \frac{N}{16Td}\left[N(d-2)-2d\right]
  = -\frac{N}{8T}
    -\frac{N^2}{16T}\epsilon
    -\frac{N^2}{16T}\epsilon^2
    +\mathcal O(\epsilon^3).
  \label{eq:evanescent-coefficient-expansion}
\end{equation}
Combining Eqs.~\eqref{eq:massive-tadpole-dimreg} and
\eqref{eq:evanescent-coefficient-expansion}, the bare figure-eight graph through
finite order is
\begin{align}
  \Delta\rho_{\mathrm{fig8}}^{\mathrm{bare}}
  = -\frac{N m^4}{128\pi^2T}
  \Bigg[
    &\frac{1}{\epsilon^2}
    +\frac{1}{\epsilon}\left(2L+\frac{N}{2}\right)
    +2L^2+\frac{\pi^2}{6}
    +NL+\frac{N}{2}
    +\mathcal O(\epsilon)
  \Bigg].
  \label{eq:figure-eight-bare-expanded}
\end{align}
Equation~\eqref{eq:figure-eight-bare-expanded} is the expansion of the bare
connected vacuum graph, not by itself a scheme-independent physical vacuum
energy.  A fully renormalized result at order $1/T$ must include the allowed
local vacuum-energy counterterm and, depending on the chosen renormalization
conditions for the phenomenological mass deformation, any lower-order
counterterm insertions required to renormalize the parameters entering the
one-loop determinant.  Therefore a finite constant extracted from
Eq.~\eqref{eq:figure-eight-bare-expanded} alone should not be interpreted as a
physical prediction.

For comparison, the massive free fields also generate the Gaussian one-loop
term
\begin{equation}
  \rho_{1\text{-loop}}
  = \frac{N}{2}\int_p\log(p^2+m^2).
  \label{eq:gaussian-one-loop-vacuum}
\end{equation}
After $\overline{\mathrm{MS}}$ subtraction this can be written, up to an
arbitrary local vacuum-energy counterterm, as
\begin{equation}
  \rho_{1\text{-loop}}^{\overline{\mathrm{MS}}}
  = \frac{Nm^2}{8\pi}
    \left(1+\log\frac{\bar\mu^2}{m^2}\right).
  \label{eq:gaussian-one-loop-vacuum-msbar}
\end{equation}
Thus the perturbative ordering is: the determinant
Eq.~\eqref{eq:gaussian-one-loop-vacuum} is a one-loop effect already present in
the quadratic theory, whereas Eq.~\eqref{eq:figure-eight-compact-result} is the
first interaction correction and is a two-loop graph of order $1/T$.

Several checks are useful:
\begin{enumerate}
  \item \textit{Mass dimensions.}
  In two worldsheet dimensions, $[T]=2$ and $[m]=1$, so
  $m^4/T$ has dimension two, as required for an energy density.

  \item \textit{Large-tension limit.}
  The interaction correction vanishes as $T\to\infty$, consistent with the
  canonically normalized Nambu--Goto interactions being suppressed by $1/T$.

  \item \textit{Massless limit in dimensional regularization.}
  When $m\to0$, the massive tadpole becomes scaleless and the figure-eight
  correction vanishes.  This statement is specific to infinite-volume
  dimensional regularization and does not imply a vanishing finite-size
  Casimir correction.

  \item \textit{$N=1$ check.}
  For a single massive scalar, $B=A^2$ identically.  In $d=2$,
  Eq.~\eqref{eq:massive-only-quartic-euclidean} becomes
  $\mathcal L_{4,E}^{(r)}=-A^2/(16T)$ and
  $\langle A^2\rangle_0=2K_m^2$, giving
  $\Delta\rho=-K_m^2/(8T)$, in agreement with
  Eq.~\eqref{eq:strict-two-dimensional-coefficient} at $N=1$.

  \item \textit{Original-field normalization check.}
  Setting $X_{\mathrm s}=0$ implies $X_1=X_{\mathrm r}/2$ and
  $X_2=-X_{\mathrm r}/2$.  Substituting this directly into the two copies of
  Eq.~\eqref{eq:single-string-ng-four-derivative} gives
  $-T(P_{X_{\mathrm r}})^2/64
   +T(G^{X_{\mathrm r}}_{ab}G_{X_{\mathrm r}}^{ab})/32$,
  exactly the pure-relative part of
  Eq.~\eqref{eq:two-string-ng-sum-relative}.  Canonical normalization then
  reproduces Eq.~\eqref{eq:massive-only-quartic-euclidean} after Wick rotation.

  \item \textit{Early-$d=2$ warning.}
  The leading coefficient of the regulated diagram is $-N/(8T)$, but the
  evanescent $N(d-2)$ term in Eq.~\eqref{eq:figure-eight-general-d} must be
  retained whenever pole or finite terms are being quoted.  This is an
  especially useful check against an apparently plausible but incomplete
  calculation.
\end{enumerate}
